\chapter{Аналитическая часть}

\section{Описание задачи и выбранных инструментов}
Целью данной работы является разработка программного обеспечения для извлечения структурированных данных из документов формата PDF. Ключевая особенность решения --- использование больших языковых моделей (БЯМ) для генерации регулярных выражений и логики валидации данных. В рамках исследования были выбраны три передовые модели, представляющие различные классы архитектур и подходы к обучению: Gemini~3 Pro Preview, YandexGPT~5.1 и Qwen3 Coder.

Ниже приведено подробное описание каждой из используемых моделей, их архитектурные особенности и обоснование выбора для решения поставленной задачи.

\section{Характеристика используемых больших языковых моделей}

\subsection{Gemini 3 Pro Preview (Google)}
Gemini~3 Pro Preview --- новейшая модель от Google DeepMind, представленная в ноябре 2025 года~\cite{gemini3_release}. Она позиционируется как мультимодальная модель следующего поколения с улучшенными способностями к сложным рассуждениям.

Ключевые характеристики модели:
\begin{itemize}
    \item модель построена на базе усовершенствованной трансформерной архитектуры с глубокой оптимизацией механизма внимания, что обеспечивает высокую производительность при обработке длинных контекстов;
    \item поддержка входного контекста объёмом до 1~миллиона токенов позволяет загружать содержимое PDF-файлов или документацию библиотек целиком для анализа;
    \item внедрённый механизм <<думающих>> (thinking) процессов позволяет модели выстраивать цепочку рассуждений перед генерацией ответа, что критически важно для составления сложных регулярных выражений;
    \item согласно независимым бенчмаркам, модель демонстрирует высокую точность в задачах логического вывода и генерации синтетических данных~\cite{gemini_bench}.
\end{itemize}

В рамках данной работы Gemini~3 используется как <<эталонный>> генератор кода и регулярных выражений для проверки гипотез.

\subsection{YandexGPT 5.1 (Yandex)}
YandexGPT~5.1 (YaGPT~5.1) --- флагманская генеративная модель компании Яндекс, выпущенная в третьем квартале 2025 года~\cite{yagpt5_release}. Данная версия получила существенные улучшения в понимании контекста и следовании сложным инструкциям на русском языке.

Ключевые характеристики модели:
\begin{itemize}
    \item модель оптимизирована для решения прикладных задач, включая извлечение сущностей (NER) и структурирование текста на русском языке;
    \item уровень галлюцинаций значительно снижен по сравнению с предыдущими версиями, а точность работы с фактами повышена;
    \item оптимизация инференса позволяет обрабатывать запросы с высокой скоростью, что делает модель эффективной для потоковой обработки данных;
    \item нативная поддержка русского языка и культурного контекста полезна при работе с документацией библиотек PyPDF2 и специфическими языковыми конструкциями.
\end{itemize}

Выбор YaGPT~5.1 обусловлен необходимостью проверки качества генерации регулярных выражений для специфических русскоязычных конструкций (например, падежные окончания в сочетаниях <<рисунок 1>>, <<таблица 2>>).

\subsection{Qwen3 Coder 30B-A3B-Instruct (Alibaba Cloud)}
Qwen3 Coder (версия 30B-A3B-Instruct) --- специализированная открытая модель, построенная на архитектуре Mixture-of-Experts (MoE)~\cite{qwen3_coder}. Ключевой особенностью данной модификации является использование технологии активации лишь части параметров во время вычислений.

Ключевые характеристики модели:
\begin{itemize}
    \item архитектура Sparse MoE включает 30 миллиардов общих параметров, из которых в момент инференса активируется лишь около 2.4--3 миллиардов (Active 3B), что обеспечивает высокую скорость генерации при сохранении интеллектуальных способностей <<большой>> модели;
    \item сниженные требования к вычислительным ресурсам (по сравнению с плотными моделями аналогичного размера) позволяют эффективно разворачивать инструмент локально на потребительских видеокартах;
    \item модель специализирована на задачах программирования и демонстрирует высокую точность в генерации кода на Python;
    \item поддержка инструктивного режима (Instruct) обеспечивает точное следование сложным алгоритмическим требованиям при написании функций обработки данных.
\end{itemize}


\section{Сравнительный анализ выбранных моделей}
Для наглядности характеристики выбранных БЯМ сведены в таблицу~\ref{tab:llm_comparison}.

\begin{table}[H]
    \centering
    \caption{Сравнительная характеристика используемых БЯМ}
    \label{tab:llm_comparison}
    \begin{tabular}{|p{0.25\linewidth}|p{0.22\linewidth}|p{0.22\linewidth}|p{0.22\linewidth}|}
        \hline
        \textbf{Характеристика} & \textbf{Gemini 3 Pro} & \textbf{YandexGPT 5.1} & \textbf{Qwen3 Coder} \\
        \hline
        Модификация & Preview & General & 30B-A3B-Instruct \\
        \hline
        Разработчик & Google & Yandex & Alibaba Cloud \\
        \hline
        Тип доступа & API (Proprietary) & API (Proprietary) & Open Weights \\
        \hline
        Архитектура & Multimodal Transformer & Transformer & Sparse MoE (Active 3B) \\
        \hline
        Контекстное окно & 1M токенов & 32k--128k токенов & 128k токенов \\
        \hline
        Сильные стороны & Рассуждения, длинный контекст & Русский язык, доступность & Скорость, локальный запуск \\
        \hline
    \end{tabular}
\end{table}

\clearpage
